
\section{General valuations}
\section{Discrete valuation and Dedekind rings}
\section{Krull rings}

\begin{theorem}[Theorem 10.1]\label{mcrt-auto-theorem-10-1}\dcref{10.1}\source{matsumura-commutative-ring-theory:page-0087}\uses{}
\begin{mathematicalrecord}
Let R c R? c K be as above, let rtt be the maximal ideal ofR
and p that of R?, and suppose that R #R?. Then
      (i)pcmcRcR?andp#m.
    (ii) p is a prime ideal of R and R? = R,.
   (iii) R/p is a valuation ring of the field RI/p.
   (iv) Given any valuation ring Sof the field R?/p, let S be its inverse image
in R?. Then S is a valuation ring having the same field of fractions K as R?.
\end{mathematicalrecord}
\end{theorem}
\begin{proof}
\begin{mathematicalrecord}
(i)Ifx~pthenx-?4R?,sox-?~Randhencex~R;xisnotaunitofR,
so that xEm. Also, since R # R? we have p # m.
   (ii) We know that p c R, so that p = p n R, and this is a prime ideal of R.
Since R - p c R? - p = {units of R?} we have R, c R?, and moreover by
construction, the maximal ideal of R, is contained in the maximal ideal p of
R?. Thus by (i), R, = R?.
   (iii) Write cp:R? --* R?/p for the natural map; then for XER? - p, if XER
we have cp(x)~R/p, and if x$R we have q(x)-? = cp(x-?)cR/p,                 and
therefore R/p is a valuation ring of RI/p.
   (iv) Note that p c S and S/p = S, so that if XER? and x$S then x is
a unit of R?, and cp(x)$S Thus cp(x-?) = q(x)-?ES, and hence x-i~S.
If on the other hand XEK - R? then x-rep c S, so that we have proved
that SvS-? = K. n
   The valuation ring S in (iv) is called the composite of R? and S. According
to (iii), every valuation ring of K contained in R? is obtained as the
composite of R? and a valuation ring of RI/p.
   Quite generally, we write mR for the maximal ideal of a local ring R. If R
and S are local rings with R 1 S and mR n S = m, we say that R dominates Sj
and write R 3 S. If R >, S and R # S, we write R > S.
\end{mathematicalrecord}
\end{proof}


\begin{theorem}[Theorem 10.2]\label{mcrt-auto-theorem-10-2}\dcref{10.2}\source{matsumura-commutative-ring-theory:page-0087}\uses{}
\begin{mathematicalrecord}
Let K be a field, A c K a subring, and p a prime ideal of A.
Then there exists a valuation ring R of K satisfying
         RxA      and m,nA=p.
\end{mathematicalrecord}
\end{theorem}
\begin{proof}
\begin{mathematicalrecord}
Replacing A by A, we can assume that A is a local ring with P = mA.
Now write 9 for the set of all subrings B of K containing A and such that
l$pB. Now AEF, and if $6 c F is a subset totally ordered by inclusion
then the union of all the elements of 9 is again an element of 9, so that, by
General valuations                                                  73
?lO
Zorn?s lemma, 9 has an element R which is maximal for inclusion. Since
pR # R there is a maximal ideal m of R containing pR. Then R c R,,,E~, so
that R = R,, and R is local. Also p c m and p is a maximal ideal of A, so that
m n A = p. Thus it only remains to prove that R is a valuation ring of K. If
 xEK and x#R then since R[x]$F               we have l~-pR[x],      and there
exists a relation of the form
          1 = a, + a,x + ... + a,x? with aiEpR.
Since 1 - a, is unit of R we can modify this to get a relation
   (*) 1 = b,x + ... + b,x? with &em.
Among all such relations, choose one for which n is as small as possible.
Similarly, if x-l $ R we can find a relation
   (**) 1 = ci .x-l + ... + c,,,x-~ with ciEm,
and choose one for which m is as small as possible. If n 2 m we multiply
(**) by b,x? and subtract from (*), and obtain a relation of the form (*)
but with a strictly smaller degree n, which is a contradiction; if 11< m then
we get the same contradiction on interchanging the roles of x and ,Y- l.
Thus if x&R we must have x~~ER.          n
\end{mathematicalrecord}
\end{proof}


\begin{theorem}[Theorem 10.3]\label{mcrt-auto-theorem-10-3}\dcref{10.3}\source{matsumura-commutative-ring-theory:page-0088}\uses{}
\begin{mathematicalrecord}
A valuation ring is integrally closed.
\end{mathematicalrecord}
\end{theorem}
\begin{proof}
\begin{mathematicalrecord}
Let R be a valuation ring of a field K, and let XEK be integral over R,
sothatx?+a,x?-?     +...+a,=Owitha,~R.Ifx$Rthenx-?~m,,butthen
         1 +a,x-?   +...+a,x-?=O,
and we get 1 EmR, which is a contradiction. Hence XE R. n
\end{mathematicalrecord}
\end{proof}


\begin{theorem}[Theorem 10.4]\label{mcrt-auto-theorem-10-4}\dcref{10.4}\source{matsumura-commutative-ring-theory:page-0088}\uses{}
\begin{mathematicalrecord}
Let K be a field, A c K a subring, and let B be the integral
closure of A in K. Then B is the intersection of all the valuation rings of K
containing A.
\end{mathematicalrecord}
\end{theorem}
\begin{proof}
\begin{mathematicalrecord}
Write 8 for the intersection of all valuation rings of K containing
A, SO that by the previous theorem we have B? 3 B. To prove the opposite
inclusion it is enough to show that for any element XEK which is not inte-
gral over A there is a valuation ring of K containing A but not x. Set
x-l = y. The ideal yA[y] of A[y] does not contain 1: for if 1 = a,y +
a,y2+ .-. + any? with a,eA then x would be integral over A, contradicting
the assumption. Therefore there is a maximal ideal p of A[y] containing
Y~[Y], and by Theorem 2 there exists a valuation ring R of K such that
R=A[y]      and m,nA[y]      = p. Now y = x-l~rn~, so that x$R. w
   Let K be a field and A c K a subring. If a valuation ring R of K contains A
we SaYthat R has a centre in A, and the prime ideal mR n A of A is called the
Centre of R in A. The set of valuation rings of K having a centre in A is called
the Zariski space or the Zariski Riemann surface of K over A, and written
74        Valuarion rings

Zar (K, A). We will treat Zar (K, A) as a topological space, introducing a
topology as follows.
   For x r ,..., x,EK, set U(x, ,..., x,)= Zar(K,A[x, ,..., x,]). Then since
             U(X I,... ,x,)nU(Y,,...,y,)=Uix,,...,x,,y,,       . . . . y,),
 the collection 9 = { U(x,, . . ,x,)ln > 0 and X,EK 3 is the basis for the
 open sets of a topology on Zar (K, A). That is, we take as open sets the
 subsets of Zar (K, A) which can be written as a union of elements of 9,
 As in the case of Spec, this topology is called the Zariski topology.
\end{mathematicalrecord}
\end{proof}


\begin{theorem}[Theorem 10.5]\label{mcrt-auto-theorem-10-5}\dcref{10.5}\source{matsumura-commutative-ring-theory:page-0089}\uses{}
\begin{mathematicalrecord}
Zar (K, A) is quasi-compact.
\end{mathematicalrecord}
\end{theorem}
\begin{proof}
\begin{mathematicalrecord}
It is enough to prove that if .d is a family of closed sets of Zar (K, A)
 having the finite intersection property (that is, the intersection of any finite
 number of elements of .d is non-empty) then the intersection of all the
 elements of & is non-empty. By Zorn?s lemma there exists a maximal family
 of closed sets &? having the finite intersection property and containing d.
 Since it is then enough to show that the intersection ofall the elements of&
 is non-empty, we can take LX!?= &?. Then it is easy to see that & has the
following properties:
    (4 F ,,...,F,Ed=>F,n...nF,E~,
    (p) Zr,...,Z,, are closed sets and Z,u~~~uZ,,~.d+Z~~d                for some i;
    (y) if a closed set F contains an element of d then FE&.
 For a subset F c Zar (K, A) we write F? to denote the complement of F.
 If FEd and F? = UnUA then F = nJJ;, and moreover if u(x,, . . , x,)~ =
(Jr=, U(Xi)CE~ then by (a) ab ove one of the U(x,)? must belong to d.
Hence the intersection of all the elements of d is the same thing as the
intersection of the sets of the form U(x)c belonging to d. Set
            1- = (ydc u(y-?)?Ed.).
Now since the condition for REZar(K, A) to belong to U(y-?)?                   is that
yEmR, the intersection        of all the elements of d is equal to
            {REZar(K,A)lm,         =I r}.
Write I for the ideal of A[lJ generated by r. If 1 $1 then by Theorem 2, the
above set is non-empty, as required to prove. But if 1 EZ then there is a
finite subset (y,, . . , y} c r such that 1 EC y,A[y,, . . . , y,]; but then
 U(y; ?)?n ... n U(y; l)c = 0, which contradicts the finite intersection
property of &.         n
    When K is an algebraic function field over an algebraically closed field k
of characteristic 0 (that is, K is a finitely generated extension of k), Zariski
gave a classification of valuation rings of K containing k, and using this and
the compactness result above, he succeeded in giving an algebraic proof in
characteristic 0 of the resolution of singularities of algebraic varieties of
dimension 2 and 3. However, Hironaka?s general proof of resolution ef
General valuations                                                    75

  singularities in characteristic 0 in all dimensions was obtained by other
  methods, without the use of valuation theory.
     As we saw at the beginning of this section, the ideals of R form a totally
  ordered set under inclusion. This holds not just for ideals, but for all R-
  modules contained in K. In particular, if we set
             G=(xR[xEK       and x#O],
  then G is a totally ordered set under inclusion; we will, however, give G the
  opposite order to that given by inclusion. That is, we define 6 by
           xR ,( yRoxR       =I yR.
  Mvreover, G is an Abelian group with product (xR)-(yR) = xyR. In general,
  an Abelian group H written additively, together with a total order relation
  2 is called an ordered group if the axiom
           x2y,z>t*x+z>y+t
  holds. This axiom implies
           (1) x>O,     y>O*x+y>O,          and (2) x>y+-y>-x.
  We make an ordered set Hu {co) by adding to H an element 00 bigger
  than all the elements of H, and fix the conventions co + LX= 00 for CCEH
  and a + a = cc. A mapv:K-Hu{co}               from a field K to HU{EJ}
  is called an additive valuation or just a valuation of K if it satisfies the
  conditions
      (1) VbY) = v(x) + v(y);
   :, (2) vb + y) 3 min { 44, V(Y));
    ? (3) v(x)= 000x=0.
  If we write K* for the multiplicative       group of K then u defines a
  bmomorphism       K * -H;    the image is a subgroup of H, called the value
  BtOUps of 0. We also set
   1.
               R?= {x~Klv(x) 3 0) and m, = (XEK~V(X) > 0},
     obtaining a valuation ring R, of K with m, as its maximal ideal, and call R,
     tie valuation ring of v, and m, the valuation ideal of u. Conversely, if R is a
   ?-balGation ring of K, then the group G = {xRlx~K*} described above is an
 ??tiered group, and we obtain an additive valuation of K with value group
  .l.pby defining v:K ---+ G u {CO} by u(0) = co and v(x) = xR for XEK* (there
   &Qn0 real significance in whether or not we rewrite the multiplication         in
::&
.*_I additively); the valuation ring of v is R. The additive valuation corres-
,-f:pMng        to a valuation ring R is not quite unique, but if u and v? are
*Y?$?oadditive valuations of K with value groups H and H? and both having
-$fie
 ;.y,r, valuation ring R then there exists an order-isomorphism        cp:H -H?
 .2@eh that v? = cpv(prove this!). Thus we can think of valuation rings and
         itive valuations as being two aspects of the same thing.
76       Valuation rings

   We now give some examples of ordered groups:
   (1) the additive group of real numbers R (this is isomorphic to the
multiplicative       group of positive reals), or any subgroup of this;
   (2) the group Z of rational integers;
   (3) the direct product Z? of n copies of Z, with lexicographical order,
that is
                                           the first non-zero element of
            (a,,...,a?)<(b,,...,b,)o       b -al,...,b     -a is positive
                                             1            n n
   An ordered group G is said to be Archimedean if it is order-isomorphic to
a suitable subgroup of R. The name is explained by the following theorem:
the condition in it is known as the Archimedean axiom. (Note that our
usage is completely unrelated to that in number theory, where non-
Archimedean fields are p-adic fields, as opposed to subfields of (wand @with
the usual ?Archimedean? metrics.)
\end{mathematicalrecord}
\end{proof}


\begin{theorem}[Theorem 10.6]\label{mcrt-auto-theorem-10-6}\dcref{10.6}\source{matsumura-commutative-ring-theory:page-0091}\uses{}
\begin{mathematicalrecord}
Let G be an ordered group; then G is Archimedean if and
only if the following condition holds:
    if a, LEG with a > 0, there exists a natural number n such that na > b.
\end{mathematicalrecord}
\end{theorem}
\begin{proof}
\begin{mathematicalrecord}
The condition is obviously necessary, and we prove sufficiency. If
G = (0) then we can certainly embed G in R. Suppose that G # {O}.
Fix some 0 < XEG. For any LEG, there is a well-defined largest integer
n such that nx d y (if y 2 0 this is clear by assumption; if y < 0, let m be
the smallest integer such that - y d mx, and set n = - m). Let this be no.
Now set Y, = y - nOx and let n, be the largest integer n such that nx < 10 y,;
we have 0 d n, < 10. Set y2 = lOy, - nix and let n2 be the largest integer
n such that nx < IOy,. Continuing in the same way, we find integers no,
n1,n2,..., and set q(y) = CI,where tl is the real number given by the decimal
expression n, + O.n,n,n,. . . . Then it can easily be checked that cp:G -+ R
is order-preserving, in the sense that y < y? implies q(y) < q(y?).
   We also see that cp is injective. For this, we only need to observe that
if y < y? then there exists a natural number r such that x < lo*(y? -Y);
the details are left to the reader.
   Finally we show that cpis a group homomorphism.         For LEG, we write
n/10? with neZ to denote the number obtained by taking the first r
decimal places of q(y); the numerator n is determined by the property that
nx < 10?~ < (n + 1)x. For ~?EG, if n?x < 10?~? < (n? + 1)x then we have
           (n + n?)x d lol(y + y?) < (n + n? + 2)x,
and hence
           q(y+y?)-(n+n?).lO-?<2.10-?,
so that
           l~(Y+Y?)-~(Y)-~cp(Y?)l~4.~o-?~
and since r is arbitrary, q(Y + y?) = q(Y) + &?).     B
General valuations                                                                      77
     ?lO
       A non-zero group G order-isomorphic to a subgroup of Iwis said to have
     rank 1. The rational rank of an ordered group G of rank 1 is the maximum
     number of elements of G (viewed as a subgroup of [w) which are linearly
     independent over Z. For example, the additive group G = Z + ZJ2 c iw is
     an ordered group of rank 1 and rational rank 2.
\end{mathematicalrecord}
\end{proof}


\begin{theorem}[Theorem 10.7]\label{mcrt-auto-theorem-10-7}\dcref{10.7}\source{matsumura-commutative-ring-theory:page-0092}\uses{}
\begin{mathematicalrecord}
Let R be a valuation ring having value group G. Then G has
     rank IoR has Krull dimension 1.
\end{mathematicalrecord}
\end{theorem}
\begin{proof}
\begin{mathematicalrecord}
(a) Since G # 0, R is not a field. Supposethat p is a prime ideal ofR
     distinct from mR. Let 5EmR be such that t$p, and set v(t) = x, where u is the
     additive valuation corresponding to R. Suppose that 0 # nip, and set JI =
     u(q);then yeG and x > 0, so that nx > y for some sufficiently large natural
     number n. This meansthat <?/PIER, so that YEI] R c p; then sincep is prime
     we have REP, which is a contradiction. Therefore p = (Cl).The only prime
     ideals of R are mR and (0), which means dim R = 1.
        (CL) If 0 # VEmR then mR is the unique prime ideal containing qR, and
     hence ,/(ryR) = mR. Thus for any 5~rrt~ there exists a natural number n such
     that l?~qR. From this one seeseasily that G satisfies the Archimedean
     axiom. n
\end{mathematicalrecord}
\end{proof}


\begin{theorem}[Theorem 11.3]\label{mcrt-auto-theorem-11-3}\dcref{11.3}\source{matsumura-commutative-ring-theory:page-0095}\uses{}
\begin{mathematicalrecord}
Let R be an integral domain and I a fractional ideal of R.
Then the following conditions are equivalent:
   (1) I is invertible;
   (2) I is a projective R-module;
   (3) I is finitely generated, and for every maximal ideal P of R, the
fractional ideal I, = IR, of R, is principal.
\end{mathematicalrecord}
\end{theorem}
\begin{proof}
\begin{mathematicalrecord}
(l)-(2) If I-?I = R then there exist a@ and b&?             such that
x;aibi=     l.Thena,,..., a, generate I, sincefor any x~l we have C(xbi)ai =
D VRs and Dedekind rings                                                 81
 ?ll
  x, and xb,eR. Let F = Re, + ... + Re, be the free R-module with basis
  el,. . , e,,; we define the R-linear map cp:F -I       by cp(ei) = ai, so that cp is
  surjective. Then we defined $:I -F          by writing $i:I -R         for the map
  tii(x) = b,x, and setting I)(X) = C $;(x)ei. We then have q@(x) = x, SO that cp
  splits, and I is isomorphic to a direct summand of the free module F, and
  therefore projective.
       (2)+-(l) Every R-linear map from I to R is given by multiplication         by
  some element of K (prove this!). If we let cp:F - I be a surjective map from
  a free module F = @ Rei, by assumption there exists a splitting $:I + F
  such that cp$ = 1. Write $(x) = 1 A;( x )e,.f or x~l; then by what we have said,
  each Ai determines a b,EK such that &(x) = bix, and since for each x that
  are only finitely many i such that /$(x) # 0, we have hi = 0 for all but finitely
  many i. Letting b,, . , b, be the non-zero ones, we have c aibix = x for all
  xel, where ai = q(ei). Thus c; a,b, = 1. Moreover, since hiI = Ai c R we
  have biGI- ?, and therefore I ~? I = R.
      (l)=>(3) As we have already seen, I is finitely generated. Now if c aibi = 1
 and P is any prime ideal then at least one of aibi must be a unit of R,,
 and I, = a,R,. Hence I, is a principal fractional ideal.
      (3)*(l)    If I is finitely generated then (IP?),=(Z,))?.         Indeed, the
inclusion c holds for any ideal; for 3, if I = a, R + ... + a,R and x@Z,)) ?
 then xa,ER,, so there exist c+R - P such that xa,c,ER, so that setting
 C=C 1...c, we have (cx)a,ER for all i, which gives cxeZ-?           and x@~?),.
 From the fact that I, is principal, we get ZP.(lP)- ? = R,. Now if II- ? # R
 then we can take a maximal ideal P such that II-? c P, and then
 Zp.(Zp)- l = z,jz - ?)p c PR,, which is a contradiction.       Thus we must have
II-?=R.           n

Theorem 1 I .4. Let R be a Noetherian integral domain, and P a non-zero
prime ideal of R. If P is invertible then ht P = 1 and R, is a DVR.
Proof. If P is invertible the maximal ideal PR, of R, is principal, and
condition (3) of Theorem 2 is satisfied; thus R, is a DVR, and so dim R, = 1.

 Theorem I1 S. Let R be a normal Noetherian domain. Then we have
   (i) all the prime divisors of a non-zero principal ideal have height 1;
   (ii) R = nhtp=, R,.
p?!       (i) Suppose 0 # aER and that P is one of the prime divisors of aR;
then there exists an element bE R such that aR: b = P. We set PR, = nr, and
then aR,:b=m,           so that baPIEmP?   and ba-?q!R,.  If ba-?mcm    then
bY the determinant trick ba-? is integral over R,, which contradicts the
fact that R, is integrally closed. Thus ba- ?m = R,, so that m- ?m = R,, and
then by the previous theorem we get ht m = ht P = 1.
   (ii) It is sufficient to prove that for a, bER with a ~0, bEaR, for every
82       Valuation rings

height 1 prime PESpec R implies bEaR. Let Pr,..., P, be the prime
divisors of aR, and let aR = q, n . . . nq, be a primary decomposition of
aR, where qi is a Pi-primary ideal for each i. Then since ht Pi = 1, we have
bEaR,,nR=q,      for i= l,..., n, and therefore bE fi qi = uR. H
\end{mathematicalrecord}
\end{proof}


\begin{theorem}[Theorem 11.6]\label{mcrt-auto-theorem-11-6}\dcref{11.6}\source{matsumura-commutative-ring-theory:page-0097}\uses{mcrt-auto-theorem-4-7}
\begin{mathematicalrecord}
For an integral domain R the following conditions are
equivalent:
   (1) R is a Dedekind ring;
   (2) R is either a field or a one-dimensional Noetherian normal domain;
   (3) every non-zero ideal of R can be written as a product of a finite
number of prime ideals.
    Moreover, the factorisation into primes in (3) is unique.
\end{mathematicalrecord}
\end{theorem}
\begin{proof}
\begin{mathematicalrecord}
(l)-(2)     Every non-zero ideal is invertible, and therefore finitely
generated, so that R is Noetherian. Let P be a non-zero prime ideal of R;
then by Theorem 4, the local ring R, is a DVR and ht P = 1, and therefore
either R is a field or dim R = 1. Also, by Theorem 4.7 we know that R is the
intersection of the R, as P runs through all the maximal ideals of R, but
since each R, is a DVR it follows that R is normal.
   (2) *(l) If R is a field there is no problem. If R is not a field then for every
maximal ideal P of R the local ring R, is a one-dimensional Noetherian
local ring and is normal, so that by Theorem 2 it is a principal ideal ring.
Thus by Theorem 3, R is a Dedekind ring.
   (l)*(3) Let I be a non-zero ideal. If I = R then we can view it as the
product of zero ideals; if I is itself maximal then it is the product of just one
prime ideal. We have already seen that R is Noetherian, so that we can use
descending induction on 1, that is assume that I #R and that every ideal
strictly bigger than J is a product of prime ideals. If I # R then there is a
maximal ideal P containing I, and I c IPpl c R. If IP-l = I then using
P-'P = R we would have I = IP, and by NAK this would lead to a
contradiction. Hence ZP-l #I, so that by induction we can write IP-? =
Qi.. . Q,, with Q,&pec R. Multiplying both sides by P gives I = Qi.. . Q*P.
   The proof of (3) +(l) is a little harder, and we break it up into four steps.
D VRs and Dedekind. rings                                           83
?ll

   Step 1. In general, any non-zero principal ideal aR of an integral domain
R is obviously invertible. Moreover, suppose that I and J are non-zero
fractional ideals and B = IJ; then obviously I and J invertible implies B
invertible, but the converse also holds. TO see this, from I-?J-?B    c R we
get I-?J-?   c B-? 9 and also from B-?ZJ c R we get B-?Z c J-? and
B-?J C I-?; now if B is invertible then multiplying the last two inclusions
together we get B -? =Bm?B-lIJcl-?J-l,            and hence B-? =I-?Jpl.
Therefore
         R=BB-?=IJI-?J-?=(II-?)(JJ-?),
and we must have II-?           = JJ-?   = R.

   SreP 2. Let P be a non-zero prime ideal. Let us prove that if I is an ideal
strictly bigger than P then IP = P. For this it is sufficient to show that if
I = p + aR with a$P then P c IP. Consider expressions of 1? and a2R + P
as product of prime ideals, 1? = P,. . . P, and a2R + P = Q1 . . . Q,. Then Pi
and Qj are prime ideals containing I, and so are prime ideals strictly bigger
than P, We now set R = R/P, and write - to denote the image in R of
elements or ideals of R. Then we have
          (*) P,... Pr = a2R = (il.. Q,,
and applying Step 1 to the domain R we find that pi and gj are all invertible,
and are prime ideals of 8. We can suppose that P, is a minimal element of
the set {Pl,..., P?>. Moreover, at least one of Q, ,. . . ,Q, is contained
in P,, so that we can assume that Q, c P,, and, on the other hand, since
Q1 is also a prime ideal and P r...prccQ1 we must have Q,3Bi                for
some i. Then pi c Qi c B,, and by the minimality        of Pi we have pi =
Pl = Q,. M u hi p 1 .
                   ymg through both sides of (*) by P; ? gives
        P2~..P~=Q2...Qs.
proceeding in the same way, we see that Y = s, and that after reordering the
Qi we can assume that Pi = Qi for i = 1,..., r.FromthiswegetPi=Qi,and
a2R+P=(P+aR)2        =P2 +aP+ a2R. Thus any element XGP can be
Written
         x=y+az+a?t        with YEP?, ZEP and tER.
Smce a#P we must have td,              and then as required we have
PcP2+aP=(P+aR)P.
  Step 3. Let bER      be a non-zero element; then in the factorisation
bR=p 1 . . . P,, every Pi is a maximal ideal of R. Indeed, if I is any ideal
stfictlY greater than Pi then ZPi = Pi, and by Step 1 Pi is invertible, so that
I-R.
   SteP 4. Let P be a prime ideal of R, and 0 # aGP. If aR = P,. . . P, with
PieSPec R then P must contain one of the Pi, but from Step 3 we know that
Pi is maximal, so that P = Pi. We deduce that P is a maximal ideal and is
84           Valuation rings

invertible. If every non-zero prime ideal is invertible then any non-zero ideal
of R is invertible, since it can be written as a product of primes. This
completes the proof of (3)*(l).
   Finally, if(l), (2) and (3) hold, then as we have seen in Step 2 above, the
uniqueness of factorisation into primes is a consequence of the fact that
the prime ideals of R are invertible.      n
\end{mathematicalrecord}
\end{proof}


\begin{theorem}[Theorem 11.7]\label{mcrt-auto-theorem-11-7}\dcref{11.7}\source{matsumura-commutative-ring-theory:page-0099}\uses{}
\begin{mathematicalrecord}
the Krull-Akizuki    theorem). Let A be a one-dimensional
Noetherian integral domain with field of fractions K, let L be a finite
algebraic extension held of K, and B a ring with A c B c L; then B is a
Noetherian ring of dimension at most 1, and ifJ is a non-zero ideal of B then
B/J is an A-module of finite length.
\end{mathematicalrecord}
\end{theorem}
\begin{proof}
\begin{mathematicalrecord}
We follow the method of proof of Akizuki [l] in the linear algebra
formulation of [BS]. First of all we prove the following lemma.
\end{mathematicalrecord}
\end{proof}


\begin{theorem}[Theorem 12.1]\label{mcrt-auto-theorem-12-1}\dcref{12.1}\source{matsumura-commutative-ring-theory:page-0101}\uses{}
\begin{mathematicalrecord}
If A is a Krull ring and S c A a multiplicative set, then A, is
again Krull. If F = (R},,, is a defining family of A then the subfamily
Pnln,r, where I- = (1~A1 R, 1 A,) is a defining family of A,.
\end{mathematicalrecord}
\end{theorem}
\begin{proof}
\begin{mathematicalrecord}
Setting m, for the maximal ideal of R, we have
        kreSn      m, = 0.
Let 0 # x~r)~,,R,;   there are at most finitely many REA such that Us < 0;
let A= (Al,..., 1n} be the set of these. If SEA then ,%$I-,hence we can find
t,Em, n S. Replacing tn by a suitable power, we can assumethat u,(t,x) b 0.
We then set t = nn,, t A, so that for every IeA we have u,(tx) 3 0, and
therefore txeA; but on the other hand tES so that XEA, and we have
proved that A, I nnEr R I. The opposite inclusion is obvious. The finiteness
condition (2) holds for A, so also for the subset r. n
   Krull rings defined by a finite number of DVRs have a simple structure.
Lemma 1 (Nagata). Let K be a field and R, , . . . , R, valuation rings of K;
Krull rings
 ?I*
 set A = 0 Ri. Then for any given a~ K there exists a natural number s 3 2
 such that
          (l+a+-~+a?-l)-?       and a.(1 +a+...+~?-~)-?
 both belong to A.
 pr&       We consider separately each Ri. Note first that (1 - a) (1 + a + ...
  +~-l)=l-u?.Ifu~Rithenu             -?urns, and any s >, 2 will do. If UER~ then
 Provided that there does not exist t such that 1 - u?~rn~, any s 2 2 will do.
 of I- aEmi then any s which is not a multiple of the characteristic of Ri/mi
 will do. If on the other hand 1 - u$mi but 1 - u?ern, for some t > 2, letting
 to be the smallest value of t for which this happens, we see that 1 - aSEmi
 only for s multiples of to, so that we only have to avoid these. Thus for each i
 the bad values of s (if any) are multiples of some number di > 1, so that
 choosing s not divisible by any of these di we get the result. H

 Theorem 12.2, Let K be a field and R,, . . . , R, valuation rings of K such
that Ri qkRj for i fj; set m, = rad (Ri). Then the intersection A = nl= 1Ri is
a semilocal ring, having pi = m, n A for i = 1,. . . , II as its only maximal
ideals; moreover A,, = R,. If each Ri is a DVR then A is a principal ideal
ring.
Proof. The inclusion A, c Ri is obvious. For the opposite inclusion, let
UCRi; choosing s > 2 as in the lemma, and setting u = (1 + a + ... + us- ?)- 1
we get UEA and UEA. Obviously u is a unit of Ri, so that UEA -pi
and a = (uu)/u~A,~. This proves that A,< =Ri. It follows from this
that there are no inclusions among pl, . . . , p,. If I is an ideal of A not
contained in any pi then (by Ex. 1.6) there exists xgl not contained in
u~=Ipi; then x is a unit in each Ri, and hence in A, SO that I = A. Thus
Pl,..., p, are all the maximal ideals of A.
   If each Ri is a DVR then we have mi # rn;, and hence pi # pi*?, (where p(*)
denotes p?A,n A). Thus there exists xi~pi such that x&1?, and xi$pj for
i #j; then pi = x,A. If1 is any ideal of A and ZRi = xyi Ri for i = 1,. . . , n then it
is easy to see that I = XII.. . x,yA. n
   If a Krull ring A is defined by an infinite number of DVRs then the
defining family of DVRs is not necessarily unique, but the following
theorem tells us that among them there is a minimal family.
\end{mathematicalrecord}
\end{proof}


\begin{theorem}[Theorem 12.3]\label{mcrt-auto-theorem-12-3}\dcref{12.3}\source{matsumura-commutative-ring-theory:page-0102}\uses{}
\begin{mathematicalrecord}
Let A be a Krull ring, K its field of fractions, and p a height 1
lPdmeidealofA;thenif,9=={R       1>IsA is a family of DVRs of K defining A, we
 must have A,EB. If we set 9, = {A,Ip~Spec A and htp = l} then .9-o is a
 defining family of A. Thus 9-O is the minimal defining family of DVRs of A.
 %of. By Theorem 1, A, is a KrulI ring defined by the subfamily Y1 =
 (R,IA, c R} c 9; if A, c R, then the elements of A - p are units of R,
88       Valuation rings

SO that p 3 m,nA.    If nt,n A = (0) then RI 3 K which is a contradiction,
hence m,n A # (0); since ht p = 1, we must have p = mln A. Thus if we fix
some 0 # x~p, then ul(x) > 0 for all RIeFI,    and hence F-, is a finite set,
Now by the previous theorem and Ex. 10.5, the elements of 9, correspond
bijectively with maximal ideals of A,, and 9, has just one element A,. Thus
A,EF, in other words 8, c 9.
   To prove that 9-O is a defining family of DVRs of A it is enough to show
that A 1 n,,,,=, A,. That is, it is enough to prove the implication
         for a, kA    with a # 0, beaA, for all A,ES,*~EUA.
As one sees easily, this is equivalent to saying that aA can be written as
the intersection of height 1 primary ideals. The set of REB such that
aR # R is finite, so we write R,, . . , R, for this. If we set
           aRinA = qi and rad(R,)nA = pi
then qi is a primary ideal belonging to pi for each i, and aA = ql.. . n q,.
Eliminating redundant terms from this expression, we get an irredundant
expression, say aA = q1 n. . . nq,. It is enough to show that then ht pi = 1
for 1 < i < r. By contradiction, suppose that ht p1 > 1. By Theorem 1, A,,1is
a Krull ring with defining family F? = {REBIA,,, c R), but is not itself
a DVR, and hence by Theorem 2, F? is infinite. Thus there exists R?6P-I
such that aR? = R?; we set p? = rad (R?)n A. We have a$p?, and ApI c R?
implies that p? c pl. Now by assumption aA #q, n...nq,.,         and RI is a
DVR, so that (rad(R,))? caR, for some v >O, and hence pi cql.
Therefore there exists an i 3 0 such that
           aA+p?,nq,n...nq,       and aAIpl+?nq,n...nq,
 Hence there exist by A such that b$aA but bp, c aA. In particular bp? c aA,
 but since a is a unit of R? we have
           (b/a)p? c An rad (R?) = p?.
 Taking 0 # cep? then for every n > 0 we have (b/a)?cEp? c A, and since A is
 completely integrally closed, b/aEA. This is a contradiction, and it proves
 that ht pi = 1 for 1 6 i < r. H
\end{mathematicalrecord}
\end{theorem}
\begin{proof}
\notready
\end{proof}


\begin{theorem}[Theorem 12.4]\label{mcrt-auto-theorem-12-4}\dcref{12.4}\source{matsumura-commutative-ring-theory:page-0103}\uses{mcrt-auto-theorem-11-2}
\begin{mathematicalrecord}
(i) A Noetherian normal domain is a Krull ring.
    (ii) Let A be an integral domain, K its field of fractions, and L an extension
field of K. If (Ai)i,, is a family of Krull rings contained in L satisfying
?12      Krull rings                                                           89

   the two conditions (1) A = nAi and (2) given any 0 # aeA we have
   a,-& = Ai for all but finitely many i, then A is a Krull ring.
       (iii) If A is a Krull ring then so is A[X] and A[Xj.
   ,J+oof. (i) 7%?IS f o 11ows from Theorem 11.5 and the fact that for any non-
   zero UEA there are only finitely many height 1 prime ideals containing aA
   (because these are the prime divisors of aA).
       (ii) is easy, and we leave it to the reader.
       (iii) K [X] .1s a p rincipal ideal ring and therefore a Krull ring. Moreover,
    if we let 9 be the set of height 1 prime ideals of A then for YES
   the ideal p[X] is prime in A[X], and by Theorem 11.2, (3), the local ring
   AL-Xl,,x, is a DVR of K(X). (If we write u for the additive valuation of K
   corresponding       to the valuation ring A,, we can extend u to an additive
   valuation of K(X) by setting u(F(X)) = min (u(q)} for a polynomial
            F(X) = a, f u,X + ... + u,X? (with u,EK),
   and v(F/G) = u(F) - v(G) for a rational function F(X)/G(X); then the
   ,valuation ring of u in K(X) is A [XI,,,,.)    Now we have K[X] A
   ACW,,x, = A,[X] (prove this!), and so


   by (ii) this is a Krull ring.
      Now for A[Xl], let CR},,,, be a family of DVRs of K defining A; then
   inside Ic[yXj we have A[X]l=      nnRn[X$      also by Ex. 9.5, RA[Xj is an
   integrally closed Noetherian ring, and is therefore a Krull ring by (i).
   However, we cannot use (ii) as it stands, since X is a non-unit of all the rings
   Z&[Xj, so we set R,[XI][X-I]         =Bb and note that A[Xj =K[X]n
   (n$J; now the hypothesis in (ii) is easily verified. Indeed,
             ~~(X)=U,X?+U,+,X?~+~..EA~XI]            with a,#0
   is a non-unit of B, if and only if a, is a non-unit of R,, and there are only
   finitely many such j?. Therefore A[Xj is a Krull ring. n

   Remark 2. Note that the field of fractions of A[Xl        is in general smaller
   than the field of fractions of K[TXJ.

   Remark 2. The B, occurring above are Euclidean rings ([B7], 9 1, Ex. 9).
\end{mathematicalrecord}
\end{theorem}
\begin{proof}
\notready
\end{proof}


\begin{theorem}[Theorem 12.5]\label{mcrt-auto-theorem-12-5}\dcref{12.5}\source{matsumura-commutative-ring-theory:page-0104}\uses{mcrt-auto-theorem-1-3,mcrt-auto-theorem-1-4}
\begin{mathematicalrecord}
The notions of Dedekind ring and one-dimensional Krull
    ring coincide.
\end{mathematicalrecord}
\end{theorem}
\begin{proof}
\begin{mathematicalrecord}
A Dedekind ring is a normal Noetherian domain, and therefore a
    KruII ring. Conversely, if A is a one-dimensional Krull ring, let US prove that
.r -4 is Noetherian. Let I be a non-zero ideal of A, and let 0 # u~l. If we can
?- Prove that AIaA is Noetherian then IluA is finitely generated, and thus so is
;:.
a
90       Valuation rings

1. By the corollary of Theorem 3 we can write aA = q, n.*.nq,,    where
qi are symbolic powers of prime ideals pi and pi # pj if i #j; now since
dim A = 1 each pi is maximal and we have
         A/aA = A/q, x .? x A/q,
by Theorem 1.3 and Theorem 1.4. But A/q, is a local ring with maximal
ideal pi/qi, and hence A/qi % A,,/qi A,,; now since each A, is a DVR, A/aA is
Noetherian (in fact even Artinian). Hence A is one-dimensional Noetherian
integral domain, and is normal, and is therefore a Dedekind ring. w
\end{mathematicalrecord}
\end{proof}


\begin{theorem}[Theorem 12.7]\label{mcrt-auto-theorem-12-7}\dcref{12.7}\source{matsumura-commutative-ring-theory:page-0105}\uses{mcrt-auto-theorem-3-7}
\begin{mathematicalrecord}
Y. Mori and J. Nishimura).       Let A and 9? be as in the
previous theorem. If A/p is Noetherian        for every pe.9 then A is
Noetherian.
\end{mathematicalrecord}
\end{theorem}
\begin{proof}
\begin{mathematicalrecord}
(J. Nishimura). As in the proof of Theorem 5, it is enough to show
that A/p??? is Noetherian for p ~9 and any n > 0. If n = 1 this holds by
hypothesis. For n > 1 we proceed as follows. Applying the previous
theorem with r = 1 and e = - 1 we can find an element x in the field Of
fractions of A such that u,(x) = 1 and u,(x) < 0 for every other qE9. Set
B= .4[x]. If yip then y/x~A so that y~xxB, and conversely Bc A, and
xBcpA,, so that p=xBnA.       Moreover, B=A+xB        and so
         BfxB = A/p.
Krull rings                                                                   91
ii12
Now x?B/x?+ ? B ?v B/xB for each i, so that by induction on i we see that
B/x?B is a Noetherian B-module for each i, and hence a Noetherian ring.
NOW we have


and B/x?B is a finite A/(x?B n A)-module, being generated by the imagesof
l,X,...,P,     so that by the Eakin-Nagata theorem (Theorem 3.7),
A/(~?B~A)    is Noetherian ring; therefore its quotient A/p(?) is also
Noetherian. W
\end{mathematicalrecord}
\end{proof}


\begin{theorem}[Theorem 11.2]\label{mcrt-auto-theorem-11-2}\dcref{11.2}\source{matsumura-commutative-ring-theory:page-0278}\uses{mcrt-auto-theorem-8-15,mcrt-auto-theorem-8-4,mcrt-auto-theorem-8-7}
\begin{mathematicalrecord}
3), and their fields of fractions are L and K respectively. Let
K? and K be their residue fields; then [K?:K] 6 [L:K] by Ex.10.8. Since A?/Q
is contained in the integral closure of A/P = k[X,, . . . ,X,1 in K?, by the
induction hypothesis, A?/Q is finite over A/P. Since
        Qi/Qi+l   =yliA!/yli+l
                                 A? N A?/Q and Qq= PA?

 we see that Al/PA? is finite over A/P. Moreover, A? is separated in the P-adic
topology (which is the same as the Y,-adic topology), because
 A? c kllq[Yl,. . . , ml. Since A is P-adically complete, A? is finite over A by
?Theorem 8.4. n

   From this result it is easy to derive the following theorem: Zf A is a
Noetherian local ring whose completion A* is reduced, then the integral
closure A? of A in its total ring of j-actions is finite ouer A. See [M] p. 237.
  On the other hand, if A is not reduced and if the maximal ideal m contains
a regular element (that is, a non-zero-divisor), then A? is not finite over A
(Krull [2]). In fact, if x # 0 is nilpotent, take a regular element s such that
 x$sA; (it is possible to find such an s since 0 mi = (0)). Then the elements
x/sj (forj = 1,2,3,. . . ) belong to A?. If A? is finite over A then there must be
some integer n such that s?(x/sj)~- A for all j. But then XE n,., Os?A = (0), a
contradiction.      n
    Suppose (A, m) is a one-dimensional Noetherian local integral domain.
Then A* is reduced if and only if A? is finite over A (Krull[2]).     In fact, if A? is
finite over A then it is a semilocal ring, and for each maximal ideal P of A?
the local ring(A?), is a DVR. The rad(A?)-adic topology of A? coincides with
 the m-adic topology, and the completion A?* of A? with respect to this
topology is a direct product of complete DVRs by Theorem 8.15, hence is
reduced. On the other hand it coincides with A* BAA? by Theorem 8.7.
 Since O+A-+A?         is exact, O--+A* --+ A* @ A? is also exact. Therefore
 A* is reduced. The converse holds, as already mentioned, without the
restriction on dimension.         n
    Rees [9] proved that a reduced Noetherian local ring A has reduced
 completion A* if and only if for every finite subset T of the total ring of
fractions K of A, the integral closure of A[lJ in K is finite over A[r].
    Akizuki Cl] constructed the first example of a one-dimensional Noeth-
erian local integral domain with non-reduced completion (see also Larfeldt-
Lech [ 11). To avoid such pathology, Nagata [Nl] defined and studied the
class of pseudo-geometric rings, which were called ?anneaux universelle-
ment japonais? by Grothendieck ([Gl], [G2]). These are now known as
?Nagata rings? ([Ml, [B9]). A Noetherian ring A is called a Nagata ring if
for every prime ideal P of A and for every finite extension field L of the field
 of fractions rc(P) of A/P, the integral closure of A/P in L is finite over A/P.
 For the basic properties of Nagata rings, see [M], $31. An alternative
definition is the following: a Noetherian ring A is a Nagata ring if (1) for
every maximal ideal m the formal tibres of A, are geometrically reduced,
 and (2) for every finite A-algebra B which is an integral domain, the set
 Nor(B) = {PESpecBIB,         is normal) is open in SpecB. The equivalence of
 these two definitions can easily be proved from [G2], (7.6.4) and (7.7.2).
    Although the integral closure A? of a Noetherian integral domain A may
fail to be finite over A, it is a Krull ring by the theorem of MoriCNagata
mentioned in 5 12. Because of its importance we quote here the theorem in
full.
Mori-Nagata       integral closure theorem. Let A be a Noetherian integral
domain and let A? be its derived normal ring. Then (1) A? is a Krull ring, and
(2) for every prime ideal P of A there are only finitely many prime ideals P? of
A? lying over P, and for each such P? the field of fractions rc(P?) of A//P? is
finite over K(P).

For a proof, see [Nl], (33.10). This proof depends on I.S. Cohen?s structure
\end{mathematicalrecord}
\end{theorem}
\begin{proof}
\notready
\end{proof}


\begin{theorem}[Theorem 12.6]\label{mcrt-auto-theorem-12-6}\dcref{12.6}\source{matsumura-commutative-ring-theory:page-0105}\uses{}
\begin{mathematicalrecord}
Let A be a Krull ring, K its field of fractions, and write 9
for the set of height I prime ideals of A. Suppose given any pl,. . . p,.e~
and e , , . . , ~,EZ. Then there exists XEK satisfying
         q(x) = e, for               1 <i < r
and
         u,(x)>0           for all     PEP-{pl,...,pr}.

Here ui and u,, stand for the normalised additive valuations of K
corresponding to pi and p.
\end{mathematicalrecord}
\end{theorem}
\begin{proof}
\begin{mathematicalrecord}
If y,eA is chosen so that y,ep, but y,$p?:?up,   u??upr,     then
o&i) = 6,, for 1 <if r. Similarly we choose y,,...,y,~A        such that
Ui(yj) = bij. Then we set

         y=    fi   yp';
              i=l

let pi,. . . , P: be all the primes PEP - (p,,...,pr}   for which o,(y) ~0.
Then choosing for each j = 1,. . . , s an element tjep; not belonging to
PI U?? u P,, and taking v to be sufficiently large, we see that
          x = y(t1 . . t,)?
satisfies the requirements of the theorem.     n
\end{mathematicalrecord}
\end{proof}
